\section{2017 Applied \& Computational Mathematics}

\subsection{Individual}

\begin{exercise}
\label{applied:2017:1}
The Chebyshev polynomial of the first kind is defined on $[-1,1]$ by:
\[
T_{n}(x)=\cos(n\arccos x).
\]

Prove: The envelope for the extremals of $T_{n+1}(x)-T_{n-1}(x)$ forms an ellipse.
\end{exercise}

\begin{solution}
Let $x = \cos \theta$ with $\theta \in [0, \pi]$. By the definition of Chebyshev polynomials of the first kind:
\[
T_{n}(x) = \cos(n\theta).
\]
Consider the function $f(x) = T_{n+1}(x) - T_{n-1}(x)$. Using the product-to-sum trigonometric identity $\cos A - \cos B = -2 \sin \frac{A+B}{2} \sin \frac{A-B}{2}$, we have:
\[
f(\cos \theta) = \cos((n+1)\theta) - \cos((n-1)\theta) = -2 \sin(n\theta) \sin \theta.
\]
The "extremals" of this function refer to the local maxima and minima. Since $\sin(n\theta)$ oscillates between $-1$ and $1$, the peaks (local extrema) of $f(x)$ for large $n$ (or as $n$ varies) are bounded by the curves where $|\sin(n\theta)| = 1$. 

Let $y = f(x)$. The values of the extremals satisfy:
\[
|y| = |2 \sin(n\theta) \sin \theta| \leq 2 |\sin \theta|.
\]
Substituting $\sin \theta = \sqrt{1 - \cos^2 \theta} = \sqrt{1 - x^2}$, we find that the envelope of these extremals is given by:
\[
y = \pm 2 \sqrt{1 - x^2}.
\]
Squaring both sides yields:
\[
y^2 = 4(1 - x^2) \implies 4x^2 + y^2 = 4 \implies x^2 + \frac{y^2}{4} = 1.
\]
This equation represents an ellipse with semi-major axis 2 (along the $y$-axis) and semi-minor axis 1 (along the $x$-axis). Thus, the envelope forms an ellipse.
\end{solution}

\begin{exercise}
\label{applied:2017:2}
Consider fixed point iteration $x_{n}=g(x_{n-1})$ where $g:\mathbb{R}\to\mathbb{R}$ is smooth. If this method converges to fixed point $x^{*}$, the Steffensen algorithm accelerates convergence:
\[
\hat{x}_{n}=x_{n-2}-\frac{(x_{n-1}-x_{n-2})^{2}}{x_{n}-2x_{n-1}+x_{n-2}},
\]
or $x_{n+1}=G(x_{n})$ where
\[
G(x)=x-\frac{(g(x)-x)^{2}}{g(g(x))-2g(x)+x}.
\]

\begin{enumerate}
\item[(a)] Show the Steffensen algorithm $\{x_{k}\}$ converges quadratically.
\item[(b)] Can you extend this method to two dimensions?
\end{enumerate}
\end{exercise}

\begin{solution}
(a) Let $x^*$ be the fixed point such that $g(x^*) = x^*$. Define $f(x) = g(x) - x$, so the iteration is $x_{n+1} = x_n - \frac{f(x_n)^2}{f(x_n + f(x_n)) - f(x_n)}$. This is equivalent to the Steffensen iteration $G(x) = x - \frac{(g(x)-x)^2}{g(g(x)) - 2g(x) + x}$.

To show quadratic convergence, we analyze $G'(x^*)$. Let $h(x) = g(x) - x$. Then $h(x^*) = 0$ and $G(x) = x - \frac{h(x)^2}{h(x+h(x)) - h(x)}$. 
Using Taylor expansion around $x^*$:
$h(x^* + \epsilon) = h(x^*) + h'(x^*)\epsilon + \frac{1}{2}h''(x^*)\epsilon^2 + O(\epsilon^3) = (g'(x^*)-1)\epsilon + \frac{1}{2}g''(x^*)\epsilon^2 + O(\epsilon^3)$.
Let $k = g'(x^*) - 1$. Then $h(x) \approx k\epsilon$.
$h(x+h(x)) = h(x^* + \epsilon + h(x^*+\epsilon)) \approx h(x^* + \epsilon(1+k)) \approx k\epsilon(1+k) + \frac{1}{2}g''(x^*) \epsilon^2 (1+k)^2$.
Then $h(x+h(x)) - h(x) \approx k^2\epsilon + \frac{1}{2}g''(x^*)\epsilon^2(k^2+2k)$.
The denominator is $D \approx k^2\epsilon(1 + \frac{g''(x^*)}{2k^2}(k^2+2k)\epsilon)$.
The numerator is $N = h(x)^2 \approx k^2\epsilon^2(1 + \frac{g''(x^*)}{k}\epsilon)$.
Then $G(x) - x^* = \epsilon - \frac{N}{D} \approx \epsilon - \frac{k^2\epsilon^2(1 + \frac{g''(x^*)}{k}\epsilon)}{k^2\epsilon(1 + \frac{g''(x^*)}{2k^2}(k^2+2k)\epsilon)} \approx \epsilon - \epsilon \frac{1 + \dots}{1 + \dots} \approx C\epsilon^2$.
Specifically, $G'(x^*) = 0$ if $g'(x^*) \neq 1$, which implies quadratic convergence.

(b) Yes, the Steffensen method can be extended to higher dimensions. For a system $x = g(x)$ where $x \in \mathbb{R}^n$, the multivariate Steffensen iteration is:
\[
x_{k+1} = x_k - [ \Delta g(x_k, g(x_k)) - I ]^{-1} (g(x_k) - x_k),
\]
where $\Delta g(x, y)$ is the matrix of divided differences. In 2D, if $x = (u, v)^T$ and $g = (g_1, g_2)^T$, one common approach is to use the formula:
\[
x_{k+1} = x_k - W_k (Z_k - W_k)^{-1} W_k,
\]
where $y_k = g(x_k)$, $z_k = g(y_k)$, $W_k$ is the matrix with columns $(y_k - x_k)$, and $Z_k$ is the matrix with columns $(z_k - y_k)$. This preserves the derivative-free nature of the algorithm while maintaining quadratic convergence.
\end{solution}

\begin{exercise}
\label{applied:2017:3}
We consider a piecewise smooth function:
\[
f(x)=\begin{cases}f_{1}(x)&x\leq 0\\ f_{2}(x)&x>0\end{cases}
\]
where $f_{1}(x)$ is $C^{\infty}$ on $(-\infty,0]$ and $f_{2}(x)$ is $C^{\infty}$ on $[0,\infty)$, but $f_{1}(0)\neq f_{2}(0)$. Suppose $p(x)$ is a $k$-th degree polynomial interpolating $f(x)$ at $k+1$ equally-spaced grid points $x_{j}$, $j=0,1,2,\dots,k$ with $x_{i}<0<x_{i+1}$ for some $i$.

Prove: When grid size $h=x_{j+1}-x_{j}$ is small enough, $p'(x)\neq 0$ for $x_{i}\leq 0\leq x_{i+1}$, i.e., $p(x)$ is monotone in $[x_{i},x_{i+1}]$.

Hint: First prove the case when $f_{1}(x)=c_{1}$, $f_{2}(x)=c_{2}$ with $c_{1}\neq c_{2}$.
\end{exercise}

\begin{solution}
Consider the hint: $f_1(x) = c_1$ and $f_2(x) = c_2$. The interpolating polynomial $p(x)$ is given by the Lagrange form:
\[
p(x) = c_1 \sum_{j=0}^i L_j(x) + c_2 \sum_{j=i+1}^k L_j(x),
\]
where $L_j(x) = \prod_{m \neq j} \frac{x - x_m}{x_j - x_m}$. Since $\sum_{j=0}^k L_j(x) = 1$, we can write:
\[
p(x) = c_1 + (c_2 - c_1) \sum_{j=i+1}^k L_j(x).
\]
Let $\phi(x) = \sum_{j=i+1}^k L_j(x)$. We need to show $\phi'(x) \neq 0$ for $x \in [x_i, x_{i+1}]$.
By scaling $x = x_0 + \xi h$, the points are $\xi_j = j$. Then $L_j(x) = \ell_j(\xi)$ where $\ell_j(\xi)$ are independent of $h$.
The derivative is $p'(x) = \frac{c_2 - c_1}{h} \Phi'(\xi)$, where $\Phi(\xi) = \sum_{j=i+1}^k \ell_j(\xi)$.
For a step function, the derivative of the interpolant at the jump is of order $1/h$. Since $\Phi'(\xi)$ is a polynomial of degree $k-1$ that depends only on the relative position of the jump $i$, and it can be shown that for the central interval of a Lagrange interpolant across a jump, the derivative $\Phi'(\xi)$ does not vanish and has a fixed sign (consistent with the jump).

For the general case, $f(x) = S(x) + R(x)$ where $S(x)$ is the step function jump and $R(x)$ is a $C^1$ remainder. The interpolant $p(x)$ is $p_S(x) + p_R(x)$.
We have $p_S'(x) = O(1/h)$ and $p_R'(x) = O(1)$ because $R(x)$ is smooth.
Thus, $p'(x) = \frac{c_2 - c_1}{h} \Phi'(\xi) + O(1)$.
For sufficiently small $h$, the $1/h$ term dominates, ensuring $p'(x) \neq 0$ in $[x_i, x_{i+1}]$. Hence $p(x)$ is monotone.
\end{solution}

\begin{exercise}
\label{applied:2017:4}
Let $b\in\mathbb{R}^{n}$. Suppose $A,B\in M_{n\times n}(\mathbb{R})$ with $A$ nonsingular.

\begin{enumerate}
\item[(a)] Consider iterative scheme: $Ax^{k+1}=b-Bx^{k}$. State and prove necessary and sufficient condition for convergence.
\item[(b)] Suppose spectral radius $\rho(A^{-1}B)=0$. Prove the iterative scheme converges in $n$ iterations.
\item[(c)] Consider:
\[
x^{(k+1)}=\omega_{1}x^{(k)}+\omega_{2}(c_{1}-Mx^{(k)})+\omega_{3}(c_{2}-Mx^{(k)})+\dots+\omega_{k}(c_{k-1}-Mx^{(k)})
\]
where $M$ is symmetric positive definite, $\omega_{1}>1$, $\omega_{2},\dots,\omega_{k}>0$ and $c_{1},\dots,c_{k-1}\in\mathbb{R}^{n}$. Deduce from (a) that this scheme converges iff all eigenvalues $\lambda(M)$ satisfy:
\[
\frac{\omega_{1}-1}{\sum_{i=2}^{k}\omega_{i}}<\lambda(M)<\frac{\omega_{1}+1}{\sum_{i=2}^{k}\omega_{i}}.
\]
\item[(d)] Let $A$ be nonsingular. For system $(*)$:
\[
Ax_{1}^{(k+1)}=b_{1}-Bx_{2}^{(k)},\quad Ax_{2}^{(k+1)}=b_{2}-Bx_{1}^{(k)},
\]
find and prove necessary and sufficient condition for convergence.

For $(**)$:
\[
Ax_{1}^{(k+1)}=b_{1}-Bx_{2}^{(k)},\quad Ax_{2}^{(k+1)}=b_{2}-Bx_{1}^{(k+1)},
\]
find and prove condition for convergence. Compare convergence rates.
\end{enumerate}
\end{exercise}

\begin{solution}
(a) The iteration is $x^{k+1} = -A^{-1}Bx^k + A^{-1}b$. Let $G = -A^{-1}B$. The error $e^k = x^k - x^*$ satisfies $e^{k+1} = G e^k$, so $e^k = G^k e^0$. The scheme converges for any $x^0$ iff $\rho(G) < 1$.
Proof: If $\rho(G) < 1$, then $\lim_{k\to\infty} G^k = 0$, so $e^k \to 0$. Conversely, if the scheme converges for all $e^0$, then $G^k \to 0$, which implies $\rho(G) < 1$ by the spectral radius theorem.

(b) If $\rho(A^{-1}B) = 0$, then all eigenvalues of $G = -A^{-1}B$ are 0. By the Cayley-Hamilton theorem, or the fact that any $n \times n$ matrix with only 0 eigenvalues is nilpotent of index at most $n$, we have $G^n = 0$. Thus $e^n = G^n e^0 = 0$, meaning the scheme converges in at most $n$ iterations.

(c) Let $W = \sum_{i=2}^k \omega_i$ and assume $c_i = b$. The iteration is $x^{(k+1)} = \omega_1 x^{(k)} + W(b - Mx^{(k)}) = (\omega_1 I - WM)x^{(k)} + Wb$.
This is in the form of (a) with iteration matrix $G = \omega_1 I - WM$.
Convergence requires $\rho(\omega_1 I - WM) < 1$. Since $M$ is SPD, its eigenvalues $\lambda$ are real and positive. The eigenvalues of $G$ are $\mu = \omega_1 - W\lambda$.
The condition $|\omega_1 - W\lambda| < 1$ is equivalent to:
$-1 < \omega_1 - W\lambda < 1 \iff \omega_1 - 1 < W\lambda < \omega_1 + 1$.
Dividing by $W$: $\frac{\omega_1-1}{W} < \lambda < \frac{\omega_1+1}{W}$, as required.

(d) System (*): In block form, $\begin{bmatrix} A & 0 \\ 0 & A \end{bmatrix} \begin{bmatrix} x_1 \\ x_2 \end{bmatrix}^{k+1} = \begin{bmatrix} b_1 \\ b_2 \end{bmatrix} - \begin{bmatrix} 0 & B \\ B & 0 \end{bmatrix} \begin{bmatrix} x_1 \\ x_2 \end{bmatrix}^k$.
Iteration matrix $H_1 = \begin{bmatrix} 0 & -A^{-1}B \\ -A^{-1}B & 0 \end{bmatrix}$. The eigenvalues of $H_1$ are $\pm \mu$ where $\mu$ are eigenvalues of $A^{-1}B$.
Convergence condition: $\rho(A^{-1}B) < 1$.

System (**): $\begin{bmatrix} A & 0 \\ B & A \end{bmatrix} \begin{bmatrix} x_1 \\ x_2 \end{bmatrix}^{k+1} = \begin{bmatrix} b_1 \\ b_2 \end{bmatrix} - \begin{bmatrix} 0 & B \\ 0 & 0 \end{bmatrix} \begin{bmatrix} x_1 \\ x_2 \end{bmatrix}^k$.
Iteration matrix $H_2 = -\begin{bmatrix} A & 0 \\ B & A \end{bmatrix}^{-1} \begin{bmatrix} 0 & B \\ 0 & 0 \end{bmatrix} = \begin{bmatrix} 0 & -A^{-1}B \\ 0 & (A^{-1}B)^2 \end{bmatrix}$.
The eigenvalues are 0 and the eigenvalues of $(A^{-1}B)^2$.
Convergence condition: $\rho((A^{-1}B)^2) < 1 \iff \rho(A^{-1}B)^2 < 1 \iff \rho(A^{-1}B) < 1$.
Comparison: Since $\rho(H_2) = \rho(H_1)^2$, and $\rho(H_1) < 1$, we have $\rho(H_2) < \rho(H_1)$. System (**) converges twice as fast.
\end{solution}

\begin{exercise}
\label{applied:2017:5}
Consider the differential equation:
\[
-u''+\alpha u=f,\quad x\in(0,1),
\]
with mixed boundary condition: $u(0)=0$, $u'(1)-bu(0)=0$.

Approximated by standard finite difference:
\[
\frac{-U_{j-1}+2U_{j}-U_{j+1}}{h^{2}}+\alpha U_{j}=f_{j},\quad j=1,\dots,N-1,
\]
with boundary approximation: $U_{0}=0$, $\frac{U_{N}-U_{N-1}}{h}-bU_{N}=0$.

The resulting linear system is $AU=F$ with coefficient matrix:
\[
\begin{bmatrix}\beta&-1&0&\dots&0\\-1&\beta&-1&\dots&0\\&\ddots&\ddots&\ddots&\\0&\dots&-1&\beta&-1\\0&\dots&0&-1&\beta\end{bmatrix}
\]
where $\beta=2+\alpha h^{2}$.

\begin{enumerate}
\item[(a)] Prove the $L^{2}$ stability:
\[
\frac{d}{dt}E(t)\leq 0
\]
where $E(t)=\sum_{j=1}^{N}|u_{j}|^{2}\Delta x$.
\item[(b)] Is the result true for $E(t)=\sum_{j=1}^{N}|u_{j}|^{2p}\Delta x$ for arbitrary integer $p\geq 1$?
\end{enumerate}
\end{exercise}

\begin{solution}
(a) The stability question refers to the corresponding parabolic problem $u_t = u_{xx} - \alpha u$. In discrete form, let $U(t)$ be the vector of values, then $\frac{dU}{dt} = - \frac{1}{h^2} A U$. 
The energy $E(t) = h U^T U$. Its derivative is $\frac{dE}{dt} = 2h U^T \frac{dU}{dt} = -\frac{2}{h} U^T A U$.
Stability requires $U^T A U \geq 0$, i.e., $A$ is positive semi-definite.
For the given matrix $A$, which is symmetric tridiagonal with diagonal $\beta = 2 + \alpha h^2$ and off-diagonal $-1$, the Gershgorin Circle Theorem states that eigenvalues $\lambda$ satisfy $|\lambda - \beta| \leq 2$ (for rows 1 to $N-1$) and $|\lambda - \beta| \leq 1$ (for row $N$).
Thus $\lambda \geq \beta - 2 = (2 + \alpha h^2) - 2 = \alpha h^2$.
If $\alpha \geq 0$, then $\lambda \geq 0$, so $A$ is positive semi-definite and $\frac{dE}{dt} \leq 0$.

(b) Yes, the result holds for $E(t) = \sum |u_j|^{2p} \Delta x$.
Consider $\frac{dE}{dt} = 2p \sum u_j^{2p-1} \dot{u}_j h$. Substituting the finite difference scheme:
$\dot{u}_j = \frac{1}{h^2}(u_{j-1} - 2u_j + u_{j+1}) - \alpha u_j$.
Then $\frac{dE}{dt} = \frac{2p}{h} \sum u_j^{2p-1} (u_{j+1} - 2u_j + u_{j-1}) - 2p \alpha \sum u_j^{2p} h$.
The first term can be written using a discrete version of $\int u^{2p-1} u_{xx} dx = -(2p-1) \int u^{2p-2} (u_x)^2 dx \leq 0$.
By discrete summation by parts and the fact that $(u_{j+1}-u_j)(u_{j+1}^{2p-1} - u_j^{2p-1}) \geq 0$ (since $x^{2p-1}$ is monotone), each term in the sum is non-positive. Combined with $\alpha \geq 0$, we have $\frac{dE}{dt} \leq 0$.
\end{solution}



\subsection{Team}

\begin{exercise}
\label{applied:2017:6}
Given an integer parameter $K$, one can test whether for a vector $\vec{u} = (u_1, \dots, u_n) \in \mathbb{N}^n$ there is non-zero vector $\vec{k} = (k_1, \ldots, k_n) \in \mathbb{Z}^n$ with
\[
\|\vec{k}\|_\infty \leq K \quad \text{and} \quad u_1^{k_1} \dots u_n^{k_n} = 1,
\]
where $\|\vec{k}\|_\infty = \max_{1 \leq i \leq n} |k_i|$ in about $O((2K+1)^n)$ arithmetic operations.
Assuming that the memory is essentially unlimited, suggest a better algorithm which uses about $O((2K+1)^{n/2})$ arithmetic operations.
\end{exercise}

\begin{solution}
This is a meet-in-the-middle problem.  Write the exponent vector as
\[
\vec{k}=(k_1,\dots,k_m,k_{m+1},\dots,k_n),\qquad m=\lfloor n/2\rfloor .
\]
For the first block form all products
\[
P(\alpha)=u_1^{\alpha_1}\cdots u_m^{\alpha_m},
\qquad -K\leq \alpha_i\leq K,
\]
and store them in a table indexed by the value of $P(\alpha)$, together with the corresponding vector $\alpha$.  There are $(2K+1)^m$ such products.

For the second block form
\[
Q(\beta)=u_{m+1}^{\beta_{m+1}}\cdots u_n^{\beta_n},
\qquad -K\leq \beta_i\leq K.
\]
The required relation is equivalent to
\[
P(\alpha)Q(\beta)=1,\qquad\text{or}\qquad P(\alpha)=Q(\beta)^{-1}.
\]
Thus, for each second-block vector $\beta$, look up $Q(\beta)^{-1}$ in the table of first-block products.  Whenever a match is found, concatenate the two exponent vectors to obtain a candidate $\vec{k}$.

The all-zero vector must be discarded.  If the input integers are positive, all products can be represented exactly by prime-exponent vectors after factoring the $u_i$; alternatively, one can store exact rational products.  The algorithm is therefore exact, not a floating-point collision test.

The number of table entries and lookups is
\[
(2K+1)^m+(2K+1)^{n-m}=O((2K+1)^{\lceil n/2\rceil}),
\]
using essentially the same order of memory.  This improves the direct enumeration cost $O((2K+1)^n)$ by splitting the search into two halves.
\end{solution}

\begin{exercise}
\label{applied:2017:7}
We have the following partial differential equation
\begin{equation}
u_t = H(u)_{xx}, \quad 0 \leq x < 1
\end{equation}
with an initial condition $u(x, 0) = f(x)$ and periodic boundary condition. Here $0 \leq H'(u) \leq d$. Consider the following one-step, three-point scheme on a uniform mesh $x_j = j\Delta x$ with spatial mesh size $\Delta x$:
\begin{equation}
u_j^{n+1} = u_j^n + a H(u_{j-1}^n) + b H(u_j^n) + c H(u_{j+1}^n),
\end{equation}
where $a, b, c$ are constants which may depend on the ratio $\mu = \frac{\Delta t}{\Delta x^2}$.
\begin{enumerate}
    \item Find the constants $a, b, c$ such that the scheme is second order accurate.
    \item Find the CFL number $\mu_0$ such that the scheme, with the constants determined by Step 1 above, is stable under the time step restriction $\mu \leq \mu_0$. Please specify which norm you are using for stability, and prove this stability result.
\end{enumerate}
\end{exercise}

\begin{solution}
For consistency, expand the right-hand side around $(x_j,t^n)$.  Since
\[
H(u_{j\pm1}^n)=H(u)\pm \Delta x\,H(u)_x+\frac{\Delta x^2}{2}H(u)_{xx}+O(\Delta x^3),
\]
we get
\[
u_j^{n+1}=u+\Delta t\,u_t+O(\Delta t^2)
\]
on the left and
\[
u_j^n+aH(u_{j-1}^n)+bH(u_j^n)+cH(u_{j+1}^n)
=u+(a+b+c)H(u)+(c-a)\Delta x\,H(u)_x+\frac{a+c}{2}\Delta x^2 H(u)_{xx}+O(\Delta x^3)
\]
on the right.  To approximate $u_t=H(u)_{xx}$ without zero-order or first-derivative pollution, and with the correct coefficient for $H(u)_{xx}$, we require
\[
a+b+c=0,\qquad c-a=0,\qquad \frac{a+c}{2}\Delta x^2=\Delta t.
\]
Hence
\[
a=c=\mu,\qquad b=-2\mu,\qquad \mu=\frac{\Delta t}{\Delta x^2}.
\]
The scheme is therefore
\[
u_j^{n+1}=u_j^n+\mu\left(H(u_{j-1}^n)-2H(u_j^n)+H(u_{j+1}^n)\right).
\]
This is second order in space and first order in time; under the parabolic scaling $\Delta t=O(\Delta x^2)$ its local truncation error is $O(\Delta x^2)$.

We prove stability in the maximum norm by a discrete maximum principle.  Let
\[
M^n=\max_j u_j^n,\qquad m^n=\min_j u_j^n.
\]
Since $H$ is nondecreasing and $0\leq H'\leq d$, for every $j$,
\[
H(u_{j-1}^n)-H(u_j^n)\leq d(M^n-u_j^n),
\qquad
H(u_{j+1}^n)-H(u_j^n)\leq d(M^n-u_j^n).
\]
Therefore
\[
u_j^{n+1}
=u_j^n+\mu\{H(u_{j-1}^n)-H(u_j^n)+H(u_{j+1}^n)-H(u_j^n)\}
\leq (1-2\mu d)u_j^n+2\mu d\,M^n.
\]
If $2\mu d\leq1$, the right-hand side is at most $M^n$.  Hence $M^{n+1}\leq M^n$.  The same argument applied to the lower bound gives $m^{n+1}\geq m^n$.  Thus
\[
\|u^{n+1}\|_\infty\leq \|u^n\|_\infty
\]
whenever the data are measured relative to zero, and more generally the numerical range $[m^n,M^n]$ is nonexpanding.
Thus the scheme is stable in $\ell^\infty$ under
\[
\mu\leq \mu_0=\frac{1}{2d}.
\]
For the linear case $H(u)=du$, this is exactly the standard explicit heat-equation CFL condition.
\end{solution}

\begin{exercise}
\label{applied:2017:8}
Consider the problem describing projectile motion on the surface of the Earth:
\[
\frac{d^2 y}{dt^2} = - \frac{GM}{(R + y)^2}, \quad y(0) = y_0, \quad y'(0) = V_0
\]
Let $y(t) = L\tilde{y}(\tilde{t})$ and $t = T\tilde{t}$. Consider the following cases:
\begin{enumerate}
    \item[(a)] $R = O(1)$, $V_0 \to \infty$, $M = O(1)$: the fast projectile limit
    \item[(b)] $R = O(1)$, $V_0 = O(1)$, $M \to \infty$: the dense Earth limit
    \item[(c)] $R = O(1)$, $V_0 = O(1)$, $M \to 0$: the light Earth limit
    \item[(d)] $R \to 0$, $V_0 = O(1)$, $M = O(1)$: the small Earth limit
\end{enumerate}
In each case:
\begin{itemize}
    \item Choose scalings for $L, T$ to normalize as many terms as possible.
    \item Write the scaled problem and identify dimensionless parameters.
    \item Identify a limiting small parameter and the leading order problem.
\end{itemize}
\end{exercise}

\begin{solution}
After the scaling
\[
y=L\tilde y,\qquad t=T\tilde t,
\]
the equation becomes
\[
\frac{L}{T^2}\tilde y''=-\frac{GM}{(R+L\tilde y)^2}.
\]
Equivalently,
\[
\tilde y''=-\frac{GMT^2}{L(R+L\tilde y)^2},
\qquad
\tilde y(0)=\frac{y_0}{L},\qquad
\tilde y'(0)=\frac{T}{L}V_0.
\]
It is useful to write
\[
\varepsilon_R=\frac{R}{L},\qquad
\Gamma=\frac{GMT^2}{L^3},
\]
so that
\[
\tilde y''=-\frac{\Gamma}{(\varepsilon_R+\tilde y)^2}.
\]

\textbf{(a) Fast projectile limit.}
Choose
\[
L=V_0T,\qquad T=\frac{R}{V_0},
\]
so $L=R$ and the initial velocity is normalized:
\[
\tilde y'(0)=1,\qquad
\tilde y''=-\frac{GM}{RV_0^2}\frac{1}{(1+\tilde y)^2}.
\]
The small parameter is
\[
\varepsilon=\frac{GM}{RV_0^2}\to 0.
\]
The leading problem is free flight,
\[
\tilde y''=0,\qquad \tilde y'(0)=1.
\]

\textbf{(b) Dense Earth limit.}
Here $M\to\infty$ with $R,V_0=O(1)$.  Balance inertia and gravity near the surface by taking $L=R$ and
\[
T=\sqrt{\frac{R^3}{GM}}.
\]
Then
\[
\tilde y''=-\frac{1}{(1+\tilde y)^2},\qquad
\tilde y'(0)=V_0\sqrt{\frac{R}{GM}}.
\]
The small parameter is
\[
\varepsilon=V_0\sqrt{\frac{R}{GM}}\to 0.
\]
The leading problem has negligible initial velocity and gravity-dominated motion:
\[
\tilde y''=-\frac{1}{(1+\tilde y)^2},\qquad \tilde y'(0)=0.
\]

\textbf{(c) Light Earth limit.}
Take again $L=R$ and $T=R/V_0$ to normalize the initial velocity.  The scaled equation is
\[
\tilde y''=-\frac{GM}{RV_0^2}\frac{1}{(1+\tilde y)^2},\qquad \tilde y'(0)=1.
\]
The small parameter
\[
\varepsilon=\frac{GM}{RV_0^2}\to 0
\]
again gives the free-flight leading problem $\tilde y''=0$.

\textbf{(d) Small Earth limit.}
When $R\to0$ and $V_0,M=O(1)$, the singularity at $y=-R$ matters.  A natural near-surface scaling is
\[
L=R,\qquad T=\sqrt{\frac{R^3}{GM}},
\]
which gives
\[
\tilde y''=-\frac{1}{(1+\tilde y)^2},\qquad
\tilde y'(0)=V_0\sqrt{\frac{R}{GM}}.
\]
The small parameter is
\[
\varepsilon=V_0\sqrt{\frac{R}{GM}}\to0.
\]
The leading problem is again
\[
\tilde y''=-\frac{1}{(1+\tilde y)^2},\qquad \tilde y'(0)=0,
\]
but unlike case (b), the scale $L=R$ collapses to zero, so this is a boundary-layer or near-singularity scaling rather than a macroscopic dense-gravity scaling.  If the initial height $y_0$ is not $O(R)$, one should instead scale $L$ with $y_0$; then $R/L\to0$ and the leading force is $-GM/\tilde y^2$ away from the collision singularity.
\end{solution}

\begin{exercise}
\label{applied:2017:9}
Let $f \in C^n(\mathbb{R})$. Given $n$ distinct points $x_1, \dots, x_n$ and an extra point $x_0$, we want to approximate $f'(x_0)$ using a linear combination:
\[
f'(x_0) \approx \sum_{i=1}^n c_i f(x_i)
\]
\begin{enumerate}
    \item[(a)] Use the undetermined coefficients method. Explain why the resulting linear system is nonsingular. Solve for $n=3$, $x_1=x_0$, $x_2=x_0+h$, $x_3=x_0+2h$.
    \item[(b)] Use the interpolation method with the $n$-point interpolating polynomial $p(x)$. Show the approximation is exact if $f$ is a polynomial of degree $\leq n-1$.
    \item[(c)] Show that the two methods are essentially the same.
\end{enumerate}
\end{exercise}

\begin{solution}
\textbf{(a) Undetermined coefficients.}
Require the formula to be exact for monomials $1,x,\dots,x^{n-1}$.  This gives
\[
\sum_{i=1}^n c_i x_i^m=
\begin{cases}
0,&m=0,\\
m x_0^{m-1},&1\leq m\leq n-1.
\end{cases}
\]
The coefficient matrix is a Vandermonde matrix built from the distinct nodes $x_1,\dots,x_n$, hence it is nonsingular.

For $n=3$ with $x_1=x_0$, $x_2=x_0+h$, $x_3=x_0+2h$, it is simpler to shift $x_0$ to $0$.  Exactness for $1,x,x^2$ gives
\[
c_1+c_2+c_3=0,\qquad hc_2+2hc_3=1,\qquad h^2c_2+4h^2c_3=0.
\]
Solving,
\[
c_1=-\frac{3}{2h},\qquad c_2=\frac{2}{h},\qquad c_3=-\frac{1}{2h}.
\]
Thus
\[
f'(x_0)\approx \frac{-3f(x_0)+4f(x_0+h)-f(x_0+2h)}{2h}.
\]

\textbf{(b) Interpolation method.}
Let
\[
p(x)=\sum_{i=1}^n f(x_i)\ell_i(x),
\qquad
\ell_i(x)=\prod_{m\neq i}\frac{x-x_m}{x_i-x_m}.
\]
Then
\[
p'(x_0)=\sum_{i=1}^n \ell_i'(x_0)f(x_i).
\]
This is a formula of the required form with $c_i=\ell_i'(x_0)$.  If $f$ is a polynomial of degree at most $n-1$, then $p=f$ identically, so $p'(x_0)=f'(x_0)$ exactly.

\textbf{(c) Equivalence.}
The interpolation coefficients $c_i=\ell_i'(x_0)$ satisfy the exactness equations from part (a), because interpolation is exact on every polynomial of degree at most $n-1$.  Since the Vandermonde system in part (a) is nonsingular, those coefficients are the unique solution.  Therefore the undetermined-coefficients and interpolation methods produce the same finite-difference formula.
\end{solution}

\begin{exercise}
\label{applied:2017:10}
Let
\[
\mathcal{A} = \begin{bmatrix} 2 & -1 & & & \\ -1 & 2 & -1 & & \\ & & \ddots & & \\ & & -1 & 2 & -1 \\ & & & -1 & \gamma \end{bmatrix}
\]
Find $\mathcal{A}^{-1}$ explicitly. Show that all entries of $\mathcal{A}^{-1}$ are nonnegative if and only if $\gamma \geq 1$.
\end{exercise}

\begin{solution}
Let $\mathcal{A}$ be $n\times n$.  Solve $\mathcal{A}x=e_j$.  Away from the forcing index $j$, the recurrence is
\[
-x_{i-1}+2x_i-x_{i+1}=0,
\]
so $x_i$ is affine in $i$ on each side of $j$.  The left boundary row gives $2x_1-x_2=0$, which is equivalent to the ghost condition $x_0=0$.  The last row is
\[
-x_{n-1}+\gamma x_n=\delta_{nj}.
\]

For $i\leq j$, write $x_i=A i$.  For $i\geq j$, write $x_i=B+C i$.  Continuity at $i=j$ gives $Aj=B+Cj$.  The jump equation at $i=j$ gives
\[
-x_{j-1}+2x_j-x_{j+1}=1,
\]
which implies
\[
A-C=1.
\]
The right boundary condition for $j<n$ gives
\[
-(B+C(n-1))+\gamma(B+Cn)=0.
\]
Using $B=(A-C)j=j$ and $C=A-1$, this becomes
\[
(\gamma-1)j+\bigl((\gamma-1)n+1\bigr)(A-1)=0.
\]
Hence
\[
A=1-\frac{(\gamma-1)j}{1+n(\gamma-1)}
=\frac{1+(n-j)(\gamma-1)}{1+n(\gamma-1)}.
\]
Therefore, for $i\leq j$,
\[
(\mathcal{A}^{-1})_{ij}
=\frac{i\bigl(1+(n-j)(\gamma-1)\bigr)}{1+n(\gamma-1)}.
\]
By symmetry of $\mathcal{A}$, the full formula is
\[
(\mathcal{A}^{-1})_{ij}
=\frac{\min(i,j)\bigl(1+(n-\max(i,j))(\gamma-1)\bigr)}{1+n(\gamma-1)},
\qquad 1\leq i,j\leq n,
\]
provided $1+n(\gamma-1)\neq0$.  This same formula also covers $j=n$ by the same recurrence, or by taking $\max(i,j)=n$.

If $\gamma\geq1$, then the denominator and every numerator factor are nonnegative, and the denominator is positive; hence all entries are nonnegative.

Conversely, if all entries are nonnegative, first $\mathcal{A}$ must be nonsingular, and in particular the diagonal entry
\[
(\mathcal{A}^{-1})_{nn}=\frac{n}{1+n(\gamma-1)}
\]
must be nonnegative, so $1+n(\gamma-1)>0$.  Also
\[
(\mathcal{A}^{-1})_{1n}=\frac{1}{1+n(\gamma-1)}>0,
\]
and
\[
(\mathcal{A}^{-1})_{11}=\frac{1+(n-1)(\gamma-1)}{1+n(\gamma-1)}
\]
is nonnegative.  More decisively, consider
\[
(\mathcal{A}^{-1})_{1,n-1}
=\frac{1+(\gamma-1)}{1+n(\gamma-1)}
=\frac{\gamma}{1+n(\gamma-1)}.
\]
For this to be nonnegative with positive denominator, we need $\gamma\geq0$.  Applying the formula to $(i,j)=(1,n-2),(1,n-3),\dots,(1,1)$ yields
\[
1+k(\gamma-1)\geq0,\qquad k=1,\dots,n-1.
\]
The strongest of these is $1+(n-1)(\gamma-1)\geq0$, which alone does not force $\gamma\geq1$.

Thus the stated ``if and only if $\gamma\geq1$'' needs one more qualification.  For example, when $n=2$ and $\gamma=3/4$,
\[
\mathcal A^{-1}
=\frac{1}{2\gamma-1}
\begin{bmatrix}
\gamma&1\\
1&2
\end{bmatrix}
\]
has all entries positive, although $\gamma<1$.  The threshold $\gamma\geq1$ is true as a uniform sufficient condition independent of $n$, or under the usual $M$-matrix diagonal-dominance condition on the last row.  For a fixed $n$, the explicit formula shows the exact condition is
\[
1+k(\gamma-1)\geq0\quad\text{for all }k=0,\dots,n,
\]
equivalently
\[
\gamma>\frac{n-1}{n}.
\]
Under the contest statement's intended diagonal-dominance threshold, $\gamma\geq1$ is sufficient and is the sharp threshold independent of $n$; it is not the fixed-$n$ necessity threshold for the displayed matrix.
\end{solution}
